\section{GPU-Native JHQ}
\label{sec:method1}
\subsection{Execution and representation}
Batched matrix multiplication applies the square orthogonal transform. IVF
routing then identifies lists to scan; query-specific lookup tables evaluate
primary distances, and only primary survivors access the residual codes.
The transform, codebooks and hierarchy follow JHQ. The execution design changes
physical layout and primary-distance evaluation while preserving selective
refinement.

Warp lanes score different candidates at the same subspace $m$. A scan copy in
subspace-major $[M,N]$ order makes their byte reads adjacent, whereas
candidate-major $[N,M]$ order makes them strided. This is standard coalescing.
Separately, packing four byte codes into a 32-bit word reduces loads, address
calculations and loop iterations, with unpacking performed in registers.
Packing changes instruction demand rather than the logical code volume.

\subsection{Cartesian-factorised primary distance tables}
\label{sec:designspace}
For subspace $m$, the primary table is
$T_m[c]=\lVert q'_m-C_m[c]\rVert_2^2$. Each codeword is a Cartesian product
of one-dimensional levels. Partitioning its eight code decisions therefore
partitions the coordinate sum exactly:
\begin{equation}
T_m[c]=T_m^{\mathrm{hi}}[c\gg4]+T_m^{\mathrm{lo}}[c\bmod16].
\label{eq:split}
\end{equation}
The two components contain 16 entries each, replacing 256 with 32. Construction
evaluates $32(D_s/2)$ coordinate contributions instead of $256D_s$: an
analytical $16\times$ reduction in that work and $8\times$ in table entries,
not a search-speedup prediction. Exactness is in real arithmetic; floating-point
association and tie order can differ. A general equal $G$-way grouping stores
$G2^{8/G}$ entries and requires $G$ lookups per candidate--subspace.

The scan admits a table into shared memory only if
\begin{equation}
40\,\texttt{BLOCK}+4MG2^{8/G}\le\texttt{smem\_optin}.
\label{eq:resid}
\end{equation}
The first term is candidate scratch. On the tested device,
$\texttt{smem\_optin}=101{,}376$ bytes leaves $94$--$59$\,KiB as
\texttt{BLOCK} rises from 128 to 1024. Full tables require $96$/$384$\,KiB
at $M=96/384$ and use the global-memory path; split tables require
$12$/$48$\,KiB and fit throughout this range. Global-memory loads may still
hit cache. This admission calculation does not measure cache misses or achieved
occupancy; the performance trade-off requires the comparisons in
Section~\ref{sec:rq3}.
\begin{figure}[tbp]
\centering
\includegraphics[width=\linewidth]{figs/fig_lut.pdf}
\caption{Cartesian factorisation and required table storage. The band is the
shared-memory budget remaining after candidate scratch; admission is distinct
from cache behaviour.}
\label{fig:lut}
\end{figure}

\subsection{Selective refinement and the policy question}
Only $c_k$ survivors are scored with their primary reconstruction plus residual.
Reading residuals during the broad scan would replace $O(c_kd)$ selective work
with work approaching $O(n_cd)$. The bounded selector's limitations are stated
in Section~\ref{sec:mono}. The remaining decision is $c_k=\alpha k$:
\emph{how many primary survivors should actually be refined}?
